\section{Variational Method Framework}\label{sec:variational}
% ============================================================
% ============================================================

% ------------------------------------------------------------
\subsection{Rayleigh--Ritz variational principle}
% ------------------------------------------------------------

The time-independent Schr\"odinger equation is
\begin{equation}\label{eq:TISE}
  \hat{H}\ket{\Psi} = E\ket{\Psi}\,.
\end{equation}
The variational principle states that for any normalized trial state $\ket{\tilde{\Psi}}$,
\begin{equation}\label{eq:variational_bound}
  \mel{\tilde{\Psi}}{\hat{H}}{\tilde{\Psi}} \geq E_0\,,
\end{equation}
with equality if and only if $\ket{\tilde{\Psi}}$ is the ground state.
More generally, the Hylleraas--Undheim--MacDonald (HUM) theorem guarantees that
for a variational calculation in an $N$-dimensional subspace, the $k$-th eigenvalue
$E_k^{(N)}$ is an upper bound to the $k$-th exact eigenvalue $E_k^{\text{exact}}$:
\begin{equation}
  E_k^{(N)} \geq E_k^{\text{exact}},\quad k = 0, 1, \ldots, N-1\,.
\end{equation}
Furthermore, enlarging the basis can only improve (lower) each variational eigenvalue:
$E_k^{(N+1)} \leq E_k^{(N)}$.

% ------------------------------------------------------------
\subsection{Expansion in a finite basis}
% ------------------------------------------------------------

We expand the trial wavefunction in a set of $N$ basis functions $\{\phi_i\}_{i=1}^{N}$:
\begin{equation}\label{eq:expansion}
  \ket{\tilde{\Psi}} = \sum_{i=1}^{N} c_i\,\ket{\phi_i}\,.
\end{equation}
The energy functional is
\begin{equation}
  E[\{c_i\}] = \frac{\mel{\tilde{\Psi}}{\hat{H}}{\tilde{\Psi}}}{\braket{\tilde{\Psi}}}
  = \frac{\sum_{i,j} c_i^* c_j\,H_{ij}}{\sum_{i,j} c_i^* c_j\,S_{ij}}\,,
\end{equation}
where $H_{ij} = \mel{\phi_i}{\hat{H}}{\phi_j}$ and $S_{ij} = \braket{\phi_i}{\phi_j}$
are the Hamiltonian and overlap matrix elements, respectively.

% ------------------------------------------------------------
\subsection{Derivation of the generalized eigenvalue problem}
% ------------------------------------------------------------

Minimizing $E$ with respect to the coefficients $c_k^*$ (treating $c_k$ and $c_k^*$ as independent):
\begin{align}
  \frac{\partial}{\partial c_k^*}\left[\sum_{i,j}c_i^*c_j\,H_{ij} - E\sum_{i,j}c_i^*c_j\,S_{ij}\right] &= 0\notag\\
  \sum_j H_{kj}\,c_j - E\sum_j S_{kj}\,c_j &= 0\,,\quad k = 1,\ldots, N\,.
\end{align}
In matrix form, this is the \emph{generalized eigenvalue problem}:
\begin{equation}\label{eq:gen_eig}
  \boxed{\mathbf{H}\,\mathbf{c} = E\,\mathbf{S}\,\mathbf{c}\,,}
\end{equation}
where $\mathbf{H}$ is the $N\times N$ Hamiltonian matrix, $\mathbf{S}$ is the $N\times N$
overlap matrix, and $\mathbf{c}$ is the coefficient vector.

\subsubsection{Standard eigenvalue problem (orthonormal basis)}

When the basis is orthonormal, $S_{ij} = \delta_{ij}$, and the generalized problem
reduces to the \emph{standard} eigenvalue problem:
\begin{equation}\label{eq:std_eig}
  \boxed{\mathbf{H}\,\mathbf{c} = E\,\mathbf{c}\,.}
\end{equation}
This is the case for the harmonic oscillator basis (Basis~I) and the DVR basis (Basis~IV).
The shifted HO basis (Basis~II) and Gaussian basis (Basis~III) have non-trivial overlap
and require either the generalized eigenvalue solver or a preliminary orthogonalization step.

% ------------------------------------------------------------
\subsection{L\"owdin canonical orthogonalization}\label{sec:Lowdin}
% ------------------------------------------------------------

When $\mathbf{S} \neq \mathbf{I}$, one can transform the generalized problem into
a standard one by the \emph{L\"owdin symmetric orthogonalization} (also called
canonical orthogonalization).

\subsubsection{Step 1: Diagonalize the overlap matrix}

Since $\mathbf{S}$ is real, symmetric, and positive semi-definite, it can be diagonalized:
\begin{equation}
  \mathbf{S} = \mathbf{U}\,\bm{\sigma}\,\mathbf{U}^T\,,
\end{equation}
where $\mathbf{U}$ is an orthogonal matrix of eigenvectors and
$\bm{\sigma} = \text{diag}(\sigma_1, \sigma_2, \ldots, \sigma_N)$
contains the eigenvalues $\sigma_i \geq 0$.

\subsubsection{Step 2: Remove near-linear dependencies}

Eigenvalues $\sigma_i < \eps_{\text{thresh}}$ (typically $\eps_{\text{thresh}} = 10^{-8}$)
indicate near-linear dependence in the basis. The corresponding eigenvectors are
discarded. Let $N' \leq N$ be the number of retained eigenvalues with
$\sigma_i \geq \eps_{\text{thresh}}$, and let $\mathbf{U}'$ and $\bm{\sigma}'$
denote the truncated matrices.

\subsubsection{Step 3: Construct the transformation matrix}

Define the matrix
\begin{equation}
  \mathbf{X} = \mathbf{U}'\,(\bm{\sigma}')^{-1/2}\,,
\end{equation}
where $(\bm{\sigma}')^{-1/2} = \text{diag}(\sigma_1^{-1/2}, \ldots, \sigma_{N'}^{-1/2})$.
This is an $N \times N'$ matrix.

\subsubsection{Step 4: Verify orthonormality}

The transformed overlap matrix is
\begin{equation}
  \mathbf{X}^T\,\mathbf{S}\,\mathbf{X}
  = (\bm{\sigma}')^{-1/2}\,(\mathbf{U}')^T\,\mathbf{U}'\,\bm{\sigma}'\,(\mathbf{U}')^T\,\mathbf{U}'\,(\bm{\sigma}')^{-1/2}
  = (\bm{\sigma}')^{-1/2}\,\bm{\sigma}'\,(\bm{\sigma}')^{-1/2}
  = \mathbf{I}_{N'}\,.
\end{equation}
The transformed basis is orthonormal.

\subsubsection{Step 5: Solve the standard eigenvalue problem}

Define the transformed Hamiltonian:
\begin{equation}
  \tilde{\mathbf{H}} = \mathbf{X}^T\,\mathbf{H}\,\mathbf{X}\,.
\end{equation}
This is an $N' \times N'$ real symmetric matrix. Diagonalize it:
\begin{equation}
  \tilde{\mathbf{H}}\,\tilde{\mathbf{c}}_k = E_k\,\tilde{\mathbf{c}}_k\,.
\end{equation}
The eigenvectors in the original (non-orthogonal) basis are recovered as
\begin{equation}\label{eq:Lowdin_backtransform}
  \mathbf{c}_k = \mathbf{X}\,\tilde{\mathbf{c}}_k\,.
\end{equation}

\begin{result}[L\"owdin procedure summary]
\begin{equation}
  \mathbf{S} \xrightarrow{\text{diag.}} \mathbf{U}\,\bm{\sigma}\,\mathbf{U}^T
  \xrightarrow{\text{truncate}} \mathbf{X} = \mathbf{U}'(\bm{\sigma}')^{-1/2}
  \xrightarrow{\text{transform}} \tilde{\mathbf{H}} = \mathbf{X}^T\mathbf{H}\mathbf{X}
  \xrightarrow{\text{diag.}} E_k,\;\mathbf{c}_k = \mathbf{X}\tilde{\mathbf{c}}_k\,.
\end{equation}
\end{result}

% ------------------------------------------------------------
\subsection{Convergence and the variational upper bound}
% ------------------------------------------------------------

By the HUM theorem, the variational eigenvalues $\{E_k^{(N)}\}$ provide upper bounds to
the exact eigenvalues and converge monotonically from above as $N \to \infty$.
A practical convergence criterion for the $k$-th level is
\begin{equation}
  |E_k^{(N)} - E_k^{(N-\delta N)}| < \delta_{\text{conv}}\,,
\end{equation}
where $\delta N$ is the basis increment and $\delta_{\text{conv}}$ is a prescribed tolerance
(e.g., $10^{-10}\,E_h$ for spectroscopic accuracy).

For the quartic double-well, the lowest levels typically converge with $N \sim 20$--$50$
harmonic oscillator functions. Higher excited states require larger bases due to the
polynomial growth of the potential.

% ============================================================
% ============================================================
